% Corrigé intégré automatiquement pour la banque Exercices.
% Source : A_Integrer/DDS_03/044_SLCI_Calculs/044_SLCI_Calculs.tex
\begin{corrigebox}[Corrigé]
\CorrigeQuestion{1}
On réduit le schéma-blocs de l'intérieur vers l'extérieur : associations en série par produit, associations en parallèle par somme et boucle de retour par $H_{BF}=\dfrac{G}{1+GH}$ (retour négatif). Après simplification, on ordonne le numérateur et le dénominateur selon les puissances de $p$, puis on identifie le gain statique, la classe et les paramètres canoniques.
\CorrigeQuestion{2}
On réduit le schéma-blocs de l'intérieur vers l'extérieur : associations en série par produit, associations en parallèle par somme et boucle de retour par $H_{BF}=\dfrac{G}{1+GH}$ (retour négatif). Après simplification, on ordonne le numérateur et le dénominateur selon les puissances de $p$, puis on identifie le gain statique, la classe et les paramètres canoniques.
\CorrigeQuestion{3}
~\\

\begin{center}
\includegraphics[width=\linewidth]{cor_01.png}
\end{center}

$\dfrac{\theta(p)}{C_r(p)}=\dfrac{1}{C_{\Omega}(p)\dfrac{K}{K_{ri}}}\dfrac{\dfrac{b}{1+\tau p}\dfrac{1}{p}}{1+\dfrac{abc}{1+\tau p}\dfrac{1}{p}}$ 
$=\dfrac{1}{ k_{1}\left( 1 + \dfrac{1}{T_{1}p} \right)\dfrac{K}{K_{ri}}}\dfrac{\dfrac{b}{1+\tau p}\dfrac{1}{p}}{1+\dfrac{abc}{1+\tau p}\dfrac{1}{p}}$

$=\dfrac{T_{1}K_{ri}p}{ k_{1}\left( T_{1}p + 1 \right)K}\cdot \dfrac{b}{p\left(1+\tau p \right)+abc}$

et $\lim\limits_{t\to\infty} \theta(t) = 1$.
\end{corrigebox}
